\documentclass[aspectratio=169,hyperref={bookmarks=false,urlcolor=blue,colorlinks=true},10pt]{beamer}
\usepackage{lmodern}
\usepackage[T1]{fontenc}
\usepackage{amsmath}
\usepackage{amssymb}
\usepackage{mathtools}
\usepackage{bm}
\usepackage{xcolor}
\usepackage[separate-uncertainty=true,multi-part-units=single,range-phrase=--,range-units=single]{siunitx}
\usepackage{url}
\usepackage{graphicx}
\usepackage{subcaption}
\usepackage{fnpct}
\usepackage{xspace}
\usepackage{listings}

\lstset{basicstyle=\ttfamily,frame=single}

\usecolortheme{default}

\setbeamertemplate{blocks}[rounded]
\setbeamertemplate{navigation symbols}{}

\setbeamertemplate{footline}{
  \hfill%
  \usebeamercolor[fg]{page number in head/foot}%
  \usebeamerfont{page number in head/foot}%
  \setbeamertemplate{page number in head/foot}[framenumber]%
  \usebeamertemplate*{page number in head/foot}\kern1em\vskip2pt%
}

\begin{document}

\begin{frame}
  \frametitle{Integers}

  Representing integers on a computer is pretty easy:
  \begin{itemize}
  \item All information is stored in sequences of bits which can either represent 0 or 1
  \item Write the integer in base 2
  \item Base 2 digits with a 0 are represented by a 0 bit, digits with a 1 are represented by a 1
    bit
  \item Reserve one bit for indicating sign of the number
    \begin{itemize}
    \item Negative numbers are a little more complicated, but still not too bad
    \end{itemize}
  \item Range of values is $[-2^{b-1},2^{b-1}-1]$ where $b$ is the number of bits available
  \end{itemize}

  \begin{table}
    \begin{tabular}{r|r}
      Decimal & Binary\\
      \hline
      0 & 0\\
      1 & 1\\
      2 & 10\\
      3 & 11\\
      4 & 100\\
      5 & 101
    \end{tabular}
  \end{table}
\end{frame}

\begin{frame}
  \frametitle{What about real numbers?}

  \begin{itemize}
  \item Inherent limitation: any range contains infinite real numbers
  \item Computers have a finite number of bits
  \item Consequence: cannot represent almost all real numbers
  \end{itemize}
\end{frame}

\begin{frame}
  \frametitle{One option: fixed point}

  Basically just use an integer but scale it down:

  \begin{table}
    \begin{tabular}{r|r}
      Value & Stored integer\\
      \hline
      0.00 & 0\\
      1.50 & 150\\
      -100.95 & -10095
    \end{tabular}
  \end{table}

  \vspace{\baselineskip}

  \begin{itemize}
  \item Works well in some scenarios
  \item Many problems for mathematical applications:
    \begin{itemize}
    \item Smaller numbers have lower precision
    \item Lots of storage needed if both small and large numbers are required
    \item Would like to be able to represent very small and very large numbers
    \end{itemize}
  \end{itemize}
\end{frame}

\begin{frame}
  \frametitle{Floating point}

  Idea: consistent precision by representing the number using scientific notation:
  \begin{align*}
    x = (-1)^s m \times 2^E,
  \end{align*}
  where
  \begin{itemize}
  \item $s$ is the \emph{sign}: $s=0$ if the number is positive, $s=1$ if the number is negative
  \item $m$ is the \emph{significand}, \emph{mantissa}, or \emph{fraction}: $1\leq m < 2$
  \item $E$ is the \emph{exponent} and is an integer
  \end{itemize}

  \vspace{\baselineskip}

  Example:
  \begin{align}
    20 = 1.25 \times 16 = (-1)^0 1.25 \times 2^4,
  \end{align}
  so $s=0$, $m=1.25$, and $E=4$.
\end{frame}

\begin{frame}
  \frametitle{Floating point consequences}

  \begin{itemize}
  \item Precision is determined by $m$, the significand
  \item Representable range is determined by $E$
  \item Details of how $m$ and $E$ are represented using bits require some thought
  \end{itemize}

  \vspace{\baselineskip}

  For an introduction to error analysis of computation with floating point numbers, see ``Numerical
  Linear Algebra'' by Trefethen \& Bau, Section III.
\end{frame}

\begin{frame}[fragile]
  \frametitle{$100+1$}

  \begin{columns}[t]
    \begin{column}{0.45\textwidth}
      What is the output?

      \begin{lstlisting}
a = 100.0f0
b = 0.0f0
for i = 1:1000
    a += 0.001f0
    b += 0.001f0
end
println(a)
println(100.0f0 + b)
println(100.0f0 +
        1000.0f0 * 0.001f0)
      \end{lstlisting}

      \pause

      Answer: $a=100.99945$, $b=100.99999$, $c=101.0$
    \end{column}
    \begin{column}{0.45\textwidth}
      Why?

      \pause

      \begin{itemize}
      \item Repeated additions accumulate error
      \item Larger numbers have lower resolution, so they accumulate more error when adding small
        numbers
      \end{itemize}
    \end{column}
  \end{columns}
\end{frame}

\begin{frame}[fragile]
  \frametitle{$5/3$}

  \begin{columns}[t]
    \begin{column}{0.45\textwidth}
      What is the output?

      \begin{lstlisting}
println((1.0f0/3.0f0)*5.0f0)
println(5.0f0/3.0f0)
      \end{lstlisting}

      \pause

      Answer: $a=1.6666667$, $b=1.6666666$
    \end{column}
    \begin{column}{0.45\textwidth}
      Why?

      \pause

      \begin{itemize}
      \item Cannot represent $1/3$ exactly
      \item Division by 3 and multiplication by $1/3$ are different
      \end{itemize}
    \end{column}
  \end{columns}
\end{frame}

\begin{frame}[fragile]
  \frametitle{$a^8$}

  \begin{columns}[t]
    \begin{column}{0.45\textwidth}
      What is the output?

      \begin{lstlisting}
a = 70.00007f0
b = a*a
b = b*b
b = b*b
println(a*a*a*a*a*a*a*a)
println(b)
println(a^8)
      \end{lstlisting}

      \pause

      Answer:
      \begin{align*}
        a &= 5.764847e14\\
        b &= 5.7648474e14\\
        c &= 5.764846e14
      \end{align*}
    \end{column}
    \begin{column}{0.45\textwidth}
      Why?

      \pause

      \begin{itemize}
      \item Floating point operations not associative in general
      \item Power function is a single operation, probably most accurate
      \item Sometimes must choose between efficiency and accuracy
      \end{itemize}
    \end{column}
  \end{columns}
\end{frame}

\begin{frame}[fragile]
  \frametitle{$x^2-y^2$}

  \begin{columns}[t]
    \begin{column}{0.45\textwidth}
      What is the output?

      \begin{lstlisting}
x = 100.0f0 + 0.001f0
y = 100.0f0 - 0.001f0

println(x*x - y*y)
println((x+y) * (x-y))
      \end{lstlisting}

      \pause

      Answer: $a=0.40039062$, $b=0.39978027$
    \end{column}
    \begin{column}{0.45\textwidth}
      Why?

      \pause

      \begin{itemize}
      \item Small errors amplified by squaring
      \item Loss of precision through cancelation error
      \item The magnification of relative error through cancelation is due to the problem being
        solved, not a property of floating point arithmetic
      \item These errors can often be avoided by restating the problem in a mathematically
        equivalent form
      \end{itemize}
    \end{column}
  \end{columns}
\end{frame}

\begin{frame}
  \frametitle{Cancelation error}

  Question: How much does an ant weigh?

  \vspace{\baselineskip}

  Experiment:
  \begin{itemize}
  \item Put the ant in my car
  \item Weigh my car
  \item Weigh the car without the ant
  \item Take the difference
  \end{itemize}
\end{frame}

\end{document}

% Local Variables:
% compile-command: "pdflatex slides.tex && pdflatex slides.tex"
% End:
